\documentclass[a4paper,12pt]{jsarticle}
\usepackage[margin=25mm]{geometry}
\usepackage{amsmath,amssymb}
\usepackage{enumitem}

\begin{document}

\begin{center}
{\LARGE 数I・A 共通テスト本番レベル（解答）}
\end{center}

\vspace{1em}

\section*{第1問}
\[
f(x)=(x-(a+1))^2+\{a^2-2a+5-(a+1)^2\}=(x-(a+1))^2+4(1-a)
\]
\begin{enumerate}[label=\arabic*.]
  \item 最小値は
  \[
  4(1-a)\quad (x=a+1 \text{で達成}).
  \]
  \item \(f(x)\le0\) が実数解をもつための必要十分条件は最小値 \(\le0\)。
  \[
  4(1-a)\le0 \iff a\ge1.
  \]
  \item 判別式
  \[
  D=\{2(a+1)\}^2-4(a^2-2a+5)=16(a-1).
  \]
  異なる2実数解 \(\iff D>0 \iff a>1\)。  
  さらに
  \[
  \alpha+\beta=2(a+1)>0,\quad \alpha\beta=a^2-2a+5=(a-1)^2+4>0
  \]
  より2解はともに正。したがって
  \[
  a>1.
  \]
  \item \(\alpha<2<\beta\) は \(f(2)<0\) と同値。
  \[
  f(2)=4-4(a+1)+a^2-2a+5=a^2-6a+5=(a-1)(a-5).
  \]
  よって
  \[
  (a-1)(a-5)<0 \iff 1<a<5.
  \]
\end{enumerate}

\section*{第2問}
\begin{enumerate}[label=\arabic*.]
  \item
  \[
  \binom{10}{4}=210.
  \]
  \item 偶数5枚、奇数5枚。和が偶数 \(\iff\) 奇数枚数が0,2,4。
  \[
  \binom{5}{0}\binom{5}{4}+\binom{5}{2}\binom{5}{2}+\binom{5}{4}\binom{5}{0}
  =5+100+5=110.
  \]
  よって
  \[
  \frac{110}{210}=\frac{11}{21}.
  \]
  \item 3の倍数は \(\{3,6,9\}\) の3枚、そうでないものは7枚。
  \[
  \binom{3}{1}\binom{7}{3}=3\cdot35=105.
  \]
  よって
  \[
  \frac{105}{210}=\frac12.
  \]
  \item 最小値を \(m\)、最大値を \(m+5\) とすると \(m=1,2,3,4,5\)。  
  その間の4個から2個選ぶので各 \(m\) について
  \[
  \binom{4}{2}=6
  \]
  通り。したがって有利事象は
  \[
  5\cdot6=30.
  \]
  よって
  \[
  \frac{30}{210}=\frac17.
  \]
\end{enumerate}

\section*{第3問}
\begin{enumerate}[label=\arabic*.]
  \item 余弦定理より
  \[
  BC^2=8^2+10^2-2\cdot8\cdot10\cos60^\circ=64+100-80=84.
  \]
  よって
  \[
  BC=2\sqrt{21}.
  \]
  \item 面積
  \[
  S=\frac12\cdot8\cdot10\sin60^\circ
  =40\cdot\frac{\sqrt3}{2}=20\sqrt3.
  \]
  \item 正弦定理 \( \dfrac{BC}{\sin A}=2R \) より
  \[
  R=\frac{2\sqrt{21}}{2\sin60^\circ}
  =\frac{2\sqrt{21}}{\sqrt3}=2\sqrt7.
  \]
  \item 高さを \(h\) とすると
  \[
  20\sqrt3=\frac12\cdot(2\sqrt{21})\cdot h
  \]
  だから
  \[
  h=\frac{20\sqrt3}{\sqrt{21}}=\frac{20}{\sqrt7}=\frac{20\sqrt7}{7}.
  \]
\end{enumerate}

\section*{第4問}
\begin{enumerate}[label=\arabic*.]
  \item 60点以上の人数は \(14+8=22\) 人。したがって相対度数は
  \[
  \frac{22}{40}=0.55.
  \]
  \item 階級値は \(10,30,50,70,90\)。平均の近似値は
  \[
  \bar{x}\approx\frac{10\cdot2+30\cdot6+50\cdot10+70\cdot14+90\cdot8}{40}
  =\frac{2400}{40}=60.
  \]
  \item
  \[
  H=50+10\frac{78-60}{18}=50+10=60.
  \]
  \item
  \[
  50+10\frac{x-60}{18}\ge65
  \iff \frac{x-60}{18}\ge1.5
  \iff x\ge87.
  \]
\end{enumerate}

\end{document}
